\section{Formazione disciplinare e multidisciplinare}

LUDI Seminar Series: ``L'importanza di fare una buona impressione: vicende e progetti editoriali della Typographia Medicea delle lingue orientali (Roma, 1584-1614)'' di Sara Fani (AMS Università di Bologna). Frequentato online, 10 ottobre 2024, 14:00-17:00.

LUDI Seminar Series: ``La biblioteconomia nel contesto internazionale. Esperienze e progetti'' di Rossana Morriello (Università di Firenze). Frequentato online, 15 novembre 2024, 10:00-12:00.

``Revealing Connections: Leveraging Linked Open Data to discover knowledge through ResearchSpace''. Seminario di una giornata su archivi semantici digitali e tecnologie Linked Open Data. Aula Giorgio Prodi, Piazza San Giovanni in Monte, 2, Bologna, 28 novembre 2024, 9:30-17:00.

``Tweets or posts on X as sources of information'' di Elisavet Chantavaridou (Università di West Attica). Frequentato online, 10 dicembre 2024, 9:00-11:00.

``nodegoat: database for historical research and the Digital Humanities'' di Pim van Bree e Geert Kessels (Lab1100). Frequentato online, 6 marzo 2025, 10:00-12:00.

``Workshop PRIN Atlas'' della Prof.ssa Marilena Daquino. Aula Affreschi, Via Zamboni 34, Bologna, 26 marzo 2025, 10:00-17:00.

``AI Literacies in the Digital Humanities: Ethics, Archives, and Critical Thinking'' di Moisés Rockembach. Frequentato online, 17 giugno 2025, 10:00-13:00.

``Biblioteche digitali e IIIF. Caratteristiche, innovazioni e applicazioni'' di Fabio Cusimano. Presentazione del libro. Frequentato online, 4 luglio 2025, 15:00-17:00.